% Ad Familiares 6.22 --- headnote
% ad-familiares-06-22, May 46 BC, Rome

\textit{Cicero to Domitius, written from Rome in May 46 BC ---
Perseus dateline \textit{Scr. Romae m. Maio a. 708 (46)}. The
correspondent is in all likelihood Cn. Domitius Ahenobarbus,
son of L. Domitius Ahenobarbus, the Pompeian consul of 54 who
fell at Pharsalus; the surviving son had been pardoned by
Caesar after the war and was now living in Italy. The
appeal to ``your mother, that excellent woman'' fits
Porcia, sister of Cato, who survived her husband and
remained in Rome (Cato himself had killed himself at Utica
the year before). The letter belongs to the post-Thapsus
group of consolatory addresses to defeated Pompeians, a
register Cicero produces almost as a genre in 46--45 BC.}

\textit{The letter is brief and severely shaped: a single
opening admission that there has been nothing to say,
followed by a long suspended period of entreaty
(``I beg and beseech you \ldots that you keep yourself safe;
that you take thought \ldots; that the lessons you have
learned \ldots you put to use; and that the loss of those
\ldots you bear''), and a closing pledge of practical
service. The tetracolon of imperatives in \S 2 is the
rhetorical centre of the letter; the closing antithesis
``if not with an even mind, then at least with a brave one''
(\textit{si non aequo animo, at forti}) gives away the
philosophical temper Cicero is working in --- the same
mode he is about to develop at length in the
\textit{Tusculan Disputations}. The pretext of the letter
is that no letter was needed; the work the letter
performs is to be present, for a man whose father and
political world have both ended.}
